\section{\texttt{giraud.py}: function-by-function implementation}
Read one-bit validation, target adapters, per-pair candidate tests, intersections and orchestration. Exact-mask knowledge is an optional stronger model. The shared SM4 engine is reused through a candidate-filter callback, not by importing hidden key material.

\subsection{Implementation topics and presentation order}
\begin{enumerate}
\item \textbf{Single-bit assumption.} Introduce \texttt{SINGLE\_BITS}, \texttt{validate\_mask} and \texttt{masks\_for}. A valid fault flips exactly one input bit. Unknown-bit mode knows the error belongs to the eight-mask set; known-bit mode receives the exact mask.
\item \textbf{AES target adapter.} Follow \texttt{aes\_fault\_encrypt}: intervene before final-round SubBytes, map a flat byte to matrix coordinates and save snapshots. No later MixColumns follows this point, so the fault reaches one output byte after ShiftRows.
\item \textbf{AES inverse-S-box test.} Explain \texttt{aes\_candidates}: for every guessed final-key byte, evaluate $S^{-1}(C_j\xor k)\xor S^{-1}(F_j\xor k)$. Require a permitted one-bit input difference, not necessarily a one-bit ciphertext XOR.
\item \textbf{Intersections and AES key assembly.} Explain \texttt{intersect\_candidates} and \texttt{recover\_byte}. Repeat observations at the same output position, assemble all sixteen unique K10 bytes in output order, and reverse the AES schedule. More than one survivor means the byte remains unresolved.
\item \textbf{SM4 single-bit adaptation.} Follow \texttt{sm4\_fault\_encrypt} and \texttt{sm4\_candidates}. Undo L before the S-box differential test, identify its affected byte, enumerate compatible inputs and convert them into subkey-byte candidates.
\item \textbf{Filter knowledge.} Explain the nested \texttt{candidate\_filter} inside \texttt{recover\_sm4\_master\_key}. It forwards the simulator mask only in known-bit mode. Unknown-bit filtering must test all eight possible bit errors even if the teaching display reveals which bit was actually injected.
\item \textbf{SM4 full recovery.} Explain reuse of the shared DFA engine: recover $rk_{31}$, undo recovered later rounds, then recover $rk_{30}$, $rk_{29}$ and $rk_{28}$. Reverse the schedule and remove FK to recover the master key. This is a Giraud-inspired SM4 model, not the unchanged AES attack.
\item \textbf{Terminal orchestration.} Follow \texttt{main}, its output-only oracle and display callbacks. Explain candidate equations, unique-byte notifications, restored views, inverse-schedule steps and final re-encryption. Compare \texttt{--known-bit}, \texttt{--step} and \texttt{--color}; logs mirror the execution to TXT and colored HTML.
\item \textbf{Verification and scope.} Check full recovery only after uniqueness and compare against the normal cipher output. Distinguish the controlled software experiment from uncharacterized physical glitches on STM32F303.
\end{enumerate}
\noindent\textbf{Suggested explanation sequence:} one-bit fault $\rightarrow$ allowed error set $\rightarrow$ inverse nonlinear test $\rightarrow$ candidate intersection $\rightarrow$ full subkey recovery $\rightarrow$ master-key verification.

\subsection{\texttt{validate\_\allowbreak{}mask}}
\begin{lightcode}
def validate_mask(mask):
\end{lightcode}
\noindent\textbf{Source:} complete \texttt{giraud.py} listing in Part 9, lines 21--24.
\begin{enumerate}
\item Accept an integer mask and test membership in SINGLE\_BITS.
\item Raise ValueError unless exactly one of the eight byte bits is set. This excludes zero, multi-bit byte errors and masks wider than one byte.
\end{enumerate}
\subsection{\texttt{masks\_\allowbreak{}for}}
\begin{lightcode}
def masks_for(first, maximum):
\end{lightcode}
\noindent\textbf{Source:} complete \texttt{giraud.py} listing in Part 9, lines 27--32.
\begin{enumerate}
\item Validate the initial mask and a maximum observation count from one through eight.
\item Put the chosen mask first, then append the other single-bit masks in their fixed order and truncate the list. Return distinct faults; repeated identical observations would not add an independent constraint.
\end{enumerate}
\subsection{\texttt{aes\_\allowbreak{}fault\_\allowbreak{}encrypt}}
\begin{lightcode}
def aes_fault_encrypt(plaintext, key, byte=0, mask=None):
\end{lightcode}
\noindent\textbf{Source:} complete \texttt{giraud.py} listing in Part 9, lines 35--58.
\begin{enumerate}
\item Validate the target byte and optional one-bit mask, then attach an observer before AES round-ten SubBytes. Translate the flat byte index to row and column and mutate that state byte only when a mask was supplied.
\item Save clean/faulted injection snapshots and subsequent final-round states, execute the real AES engine and return ciphertext plus trace. A None mask leaves a clean baseline for comparison.
\end{enumerate}
\subsection{\texttt{aes\_\allowbreak{}fault\_\allowbreak{}encrypt.observer}}
\begin{lightcode}
def observer(number, operation, state):
\end{lightcode}
\noindent\textbf{Source:} complete \texttt{giraud.py} listing in Part 9, lines 47--55.
\begin{enumerate}
\item Record the clean round-ten SubBytes input, optionally mutate one bit and record the faulted input. Record later final-round events, naming the final AddRoundKey snapshot Ciphertext.
\end{enumerate}
\subsection{\texttt{sm4\_\allowbreak{}fault\_\allowbreak{}encrypt}}
\begin{lightcode}
def sm4_fault_encrypt(plaintext, key, byte=None, mask=0, fault_round=31):
\end{lightcode}
\noindent\textbf{Source:} complete \texttt{giraud.py} listing in Part 9, lines 61--66.
\begin{enumerate}
\item When a target byte is supplied, require a valid one-bit mask.
\item Delegate to piret.sm4\_fault\_encrypt with the selected round. This reuses the same SM4 engine while constraining the experiment to the single-bit model.
\end{enumerate}
\subsection{\texttt{recover\_\allowbreak{}sm4\_\allowbreak{}master\_\allowbreak{}key}}
\begin{lightcode}
def recover_sm4_master_key(correct, fault_oracle, first_mask=1, maximum=8, known_bit=False, start_byte=0, observer=None):
\end{lightcode}
\noindent\textbf{Source:} complete \texttt{giraud.py} listing in Part 9, lines 69--77.
\begin{enumerate}
\item Build a local candidate filter that passes the exact mask onward only when known\_bit is true. Generate distinct single-bit masks with masks\_for.
\item Call the shared multi-round SM4 recovery engine with this filter, start-byte order and optional observer. Return its complete or ambiguous report without obtaining reference subkeys.
\end{enumerate}
\subsection{\texttt{recover\_\allowbreak{}sm4\_\allowbreak{}master\_\allowbreak{}key.candidate\_\allowbreak{}filter}}
\begin{lightcode}
def candidate_filter(clean, faulty, simulator_mask):
\end{lightcode}
\noindent\textbf{Source:} complete \texttt{giraud.py} listing in Part 9, lines 73--75.
\begin{enumerate}
\item Accept peeled output views and the simulator mask. Pass that mask to sm4\_candidates only in known-bit mode; pass None otherwise, deliberately preserving unknown-bit uncertainty.
\end{enumerate}
\subsection{\texttt{aes\_\allowbreak{}candidates}}
\begin{lightcode}
def aes_candidates(correct, faulty, known_mask=None):
\end{lightcode}
\noindent\textbf{Source:} complete \texttt{giraud.py} listing in Part 9, lines 80--106.
\begin{enumerate}
\item Validate both blocks and any known mask, then require exactly one changed ciphertext byte. For each K10 byte guess, remove its XOR and invert the S-box on both observed values.
\item Keep the guess if the input difference is any single-bit mask, or equals the exact known mask in the stronger model. Return the changed output index and surviving byte values.
\end{enumerate}
\subsection{\texttt{sm4\_\allowbreak{}candidates}}
\begin{lightcode}
def sm4_candidates(correct, faulty, known_mask=None):
\end{lightcode}
\noindent\textbf{Source:} complete \texttt{giraud.py} listing in Part 9, lines 109--151.
\begin{enumerate}
\item Require three unchanged output words in the current final/peeled view, invert L on the changed first word and locate the single changed S-box byte.
\item Enumerate input bytes satisfying at least one permitted one-bit differential, then XOR the observable three-word combination to obtain key-byte candidates. Return the byte index, candidate set, input set and inverse-L difference.
\end{enumerate}
\subsection{\texttt{intersect\_\allowbreak{}candidates}}
\begin{lightcode}
def intersect_candidates(previous, current):
\end{lightcode}
\noindent\textbf{Source:} complete \texttt{giraud.py} listing in Part 9, lines 154--163.
\begin{enumerate}
\item Treat None as the first observation and copy the current candidate set; otherwise intersect the previous and current sets.
\item Reject an empty result as an incompatible observation/model pair. Return all survivors, leaving ambiguity explicit instead of selecting an arbitrary byte.
\end{enumerate}
\subsection{\texttt{recover\_\allowbreak{}byte}}
\begin{lightcode}
def recover_byte(algorithm, observations):
\end{lightcode}
\noindent\textbf{Source:} complete \texttt{giraud.py} listing in Part 9, lines 166--183.
\begin{enumerate}
\item Choose aes\_candidates or sm4\_candidates from the algorithm name, then iterate supplied ciphertext observations and any optional known mask.
\item Require all observations to concern the same output/subkey byte and intersect their sets. Return the position and survivors; an empty observation sequence returns None,None and this helper does not simulate a target.
\end{enumerate}
\subsection{\texttt{main}}
\begin{lightcode}
def main(argv=None):
\end{lightcode}
\noindent\textbf{Source:} complete \texttt{giraud.py} listing in Part 9, lines 186--286.
\begin{enumerate}
\item Parse algorithm, bit-knowledge, observation and display settings and decode the target configuration. For AES, fault each input byte, locate its final output byte and intersect candidate sets until unique or exhausted.
\item Reverse K10 only after all bytes are unique; for SM4, use the output-only oracle and shared four-round recovery wrapper. Print recovered keys and verify them; return a nonzero status if any selected algorithm remains ambiguous.
\end{enumerate}
\subsection{\texttt{main.oracle}}
\begin{lightcode}
def oracle(number, byte, mask):
\end{lightcode}
\noindent\textbf{Source:} complete \texttt{giraud.py} listing in Part 9, lines 232--233.
\begin{enumerate}
\item Capture the simulation plaintext and key, request the selected SM4 round/byte/mask fault and return only the output bytes. Keep the diagnostic trace outside the solver interface.
\end{enumerate}
\subsection{\texttt{main.show}}
\begin{lightcode}
def show(event):
\end{lightcode}
\noindent\textbf{Source:} complete \texttt{giraud.py} listing in Part 9, lines 235--241.
\begin{enumerate}
\item Receive one completed SM4 candidate-history event and print raw outputs, already recovered later subkeys, peeled views and survivor counts. Optional pausing affects presentation only.
\end{enumerate}
